\begin{center}
    \vspace*{0.5cm}
    {\Huge\textbf{Estudio de métodos de Meta-Learning y Domain Adaptation}}
    
    \vspace{0.2cm}
    
    {\Large Lucio Luque Materazzi, Teo Manuel Kaucher}

    \vspace{0.3cm}
    
    {\Large lluquematerazzi@udesa.edu.ar \\ tkaucher@udesa.edu.ar}

    \vspace{0.3 cm}

    {\Large Ingeniería en Inteligencia Artificial}

    \Large I309 Visión Artificial Avanzada
    
    \vspace{0.5cm}

    Otoño 2026
    
    \vspace{0.5cm}

    \begin{abstract}
    En este trabajo se estudian y comparan distintos métodos de meta-learning y domain adaptation aplicados a tareas de clasificación de dígitos bajo escenarios de pocos ejemplos y cambio de dominio. Para ello, se utilizan los datasets MNIST, MNIST-M y SVHN, que comparten las mismas clases pero presentan diferencias visuales significativas en color, textura, fondo y contexto.
    
    En primer lugar, se evalúan modelos de few-shot meta learning, específicamente ProtoNet, MAML y ProtoMAML, entrenados de manera episódica sobre MNIST y luego analizados en dominios externos. Los resultados muestran que ProtoNet y ProtoMAML presentan mayor robustez frente al cambio de dominio que MAML, especialmente en MNIST-M y SVHN.
    
    En segundo lugar, se implementa una arquitectura DANN para estudiar si el entrenamiento adversarial mediante una Gradient Reversal Layer permite aprender representaciones más invariantes entre MNIST y MNIST-M. Los experimentos muestran que DANN mejora significativamente el rendimiento sobre el dominio objetivo respecto del baseline. 
    
    Finalmente, el análisis de los embeddings mediante t-SNE permite observar que los métodos basados en prototipos y DANN generan representaciones más alineadas entre dominios, aunque SVHN continúa siendo el conjunto más desafiante debido a su mayor diferencia visual respecto de MNIST.
    \end{abstract}

    \vspace{0.15cm}
      
\end{center}